\documentclass[10pt,a4paper]{amsart}
\usepackage[margin=19mm]{geometry}
\usepackage{amsmath,amssymb,mathtools}
\usepackage[scr=rsfs,cal=euler]{mathalfa}
\usepackage{xfrac}
\usepackage[unicode,hidelinks,bookmarks=false]{hyperref}
\newcommand{\ie}{i.e.\ }
\newcommand{\cf}{cf.\ }
\newcommand{\resp}{respectively\ }
\newcommand{\Subsection}{\subsection}
\providecommand{\nobreakdash}{\nobreak\hbox{-}}
\setlength{\emergencystretch}{2em}
\multlinegap0pt
\usepackage{tikz}
\usepackage{float}
\usepackage{amssymb}
\newtheorem{lemma}{Lemma}
\def\R{{\mathbb R}}
\title{Multiplicities of eigenvalues and quadratic representations of integers}
\author{Siqi Fu, Andrew Pendleton}
\date{Source: arXiv:2603.14748v2 (CC BY 4.0)}
\begin{document}
\maketitle

\noindent\emph{Source and licence.} Multiplicities of eigenvalues and quadratic representations of integers. Version arXiv:2603.14748v2 (CC BY 4.0), distributed by arXiv under the Creative Commons Attribution 4.0 International licence, http://creativecommons.org/licenses/by/4.0/, as recorded in the arXiv licence metadata for this exact version and shown on the abstract page https://arxiv.org/abs/2603.14748v2. Full licence notice: \texttt{licenses/arxiv-2603.14748-CC-BY-4.0.txt}.\par\smallskip

\noindent\textit{Editorial scope.} Counted unit: the article's lemma lm:2 (source0.tex lines 249--251). The article's complete original proof of it is delivered with it (source0.tex lines 253--319, one \texttt{proof} environment of 67 lines). No mathematics is written, repaired or re-derived: the statement and the proof are the article's own lines, in the article's own notation, and no retained line is substituted or re-flowed.  Retained uncounted as the source-local context the proof needs: lines 244--246.  Nothing was omitted from any retained span, so the package carries no omission record.  No retained line contains a citation, so the wrapper carries no bibliography and no bibliography entry.  No retained line contains a figure environment, an image-inclusion command or drawing code, so no image is delivered and the package declares no asset dependency.  Editorial formatting only, in addition: an AMS \texttt{amsart} preamble that restates the article's own declarations and macros used by the retained lines, namely the environments \texttt{lemma} on separate counters, together with the standard packages those lines need.  The retained environments print in this document as Lemma 1 in order of appearance, whereas the article prints the same items with the numbering of its own full document.  The complete native source of the article as distributed in the arXiv e-print for this exact version is bundled unchanged as \texttt{sources/196430/source0.tex}.
\begin{equation}\label{eq:quadratic}
\tilde f(x, y):=x^2 - 2 r \cos \theta xy + r^2 y^2= \frac{|w|^2}{|\vec{b}|^2}.
\end{equation}

\begin{lemma}\label{lm:2}
	If the Laplacian on the torus  \( T=\R^2/L \) has an eigenvalue of multiplicity \( \geq 6 \), then both \( r^2 \) and \( r\cos(\theta) \) are rational.
\end{lemma}

\begin{proof}  Under the assumption, there exist at least three pairs of integral lattice points
\[
(x_1, y_1), \; (x_2, y_2),\; \text{and}\; (x_3, y_3) \; \text{with } 0 \leq y_1 < y_2 < y_3
\]
such that
\begin{equation}\label{eq:m}
\tilde f(x_1, y_1) = \tilde f(x_2, y_2) = \tilde f(x_3, y_3).
\end{equation}
Observe that
\[
\tilde f(x,y) = \bigl( x - r\cos\theta \, y \bigr)^2 + (r\sin\theta)^2 y^2.
\]
From \eqref{eq:m} we obtain
\begin{equation}\label{eq:sin}
	\begin{aligned}
(r\sin\theta)^2 &= \frac{(x_1 - r\cos\theta \, y_1)^2 - (x_2 - r\cos\theta \, y_2)^2}{y_2^2 - y_1^2}\\
&= \frac{(x_1 - r\cos\theta \, y_1)^2 - (x_3 - r\cos\theta \, y_3)^2}{y_3^2 - y_1^2}.
\end{aligned}
\end{equation}
This leads to
\begin{align*}
	&\left[ \frac{x_1 - x_2}{y_2 - y_1} + r \cos \theta \right] \left[ \frac{x_1 + x_2}{y_2 + y_1} - r \cos \theta \right] \\
	&\quad = \left[ \frac{x_1 - x_3}{y_3 - y_1} + r \cos \theta \right] \cdot \left[ \frac{x_1 + x_3}{y_3 + y_1} - r \cos \theta \right].
\end{align*}

Set
\[
s=\frac{x_1 - x_2}{y_2 - y_1} \cdot \frac{x_1 + x_2}{y_2 + y_1}
- \frac{x_1 - x_3}{y_3 - y_1} \cdot \frac{x_1 + x_3}{y_3 + y_1}
\]
and
\[
t=\frac{x_1 + x_3}{y_3 + y_1} - \frac{x_1 - x_3}{y_3 - y_1}
- \frac{x_1 + x_2}{y_2 + y_1} + \frac{x_1 - x_2}{y_2 - y_1}.
\]
Then
\begin{equation}\label{eq:cos}
(r\cos\theta) t=s.
\end{equation}
We claim that $t\not=0$.  We prove this by contradiction. Suppose $t=0$. Then
\begin{equation}\label{eq:colinear}
	\frac{x_3y_3 - x_1y_1}{y_3^2 - y_1^2} = \frac{x_2y_2 - x_1y_1}{y_2^2 - y_1^2}.
\end{equation}
Thus, the points $(x_iy_i,\ y_i^2)$, $i=1, 2, 3$ are collinear. Let $\lambda$ be the quotient given 
by \eqref{eq:colinear} and let $\beta=x_1y_1-\lambda y_1^2$.  Then 
\[
x_i=\lambda y_i+\beta/y_i, \quad i=1, 2, 3.
\]
Since all three points $(x_i,\ y_i)$, $i=1, 2, 3$ lie on the ellipse defined by \eqref{eq:quadratic}, $y_i$, $i=1, 2, 3$ are
the solutions to  
\[
A(\lambda y+\beta/y)^2+B(\lambda y+\beta/y)y+Cy^2=D
\]
where $A=1$, $B=-2r\cos\theta$, $C=r^2$, and $D=|w|^2/|\vec{b}|^2$.  The above equation reduces to
\begin{equation}\label{eq:n2}
(A\lambda^2+B\lambda+C)y^4+(2A\lambda\beta+B\beta-D)y^2+A\beta^2=0,
\end{equation}
which is a quadratic equation in $y^2$. (Note that the coefficient of $y^4$ in the above 
equation is positive because $B^2-4AC<0$.)  It is impossible for \eqref{eq:n2} to have 
three distinct positive solutions.  This concludes the proof of the claim.
  
From \eqref{eq:cos}, we then obtain that \( r \cos \theta \in \mathbb{Q} \).  It then follows from \eqref{eq:sin} that \( (r\sin\theta)^2 \in \mathbb{Q} \).  
Therefore
\[
r^2 = (r\cos\theta)^2 + (r\sin\theta)^2 \in \mathbb{Q}.
\]
\end{proof}

\end{document}
